
\documentclass{article}
\usepackage{colortbl}
\usepackage{makecell}
\usepackage{multirow}
\usepackage{supertabular}

\begin{document}

\newcounter{utterance}

\centering \large Interaction Transcript for game `clean\_up', experiment `2\_hard\_5obj\_lv', episode 1 with mistral{-}large{-}3{-}azure.
\vspace{24pt}

{ \footnotesize  \setcounter{utterance}{1}
\setlength{\tabcolsep}{0pt}
\begin{supertabular}{c@{$\;$}|p{.15\linewidth}@{}p{.15\linewidth}p{.15\linewidth}p{.15\linewidth}p{.15\linewidth}p{.15\linewidth}}
    \# & $\;$A & \multicolumn{4}{c}{Game Master} & $\;\:$B\\
    \hline

    \theutterance \stepcounter{utterance}  
    & & \multicolumn{4}{p{0.6\linewidth}}{
        \cellcolor[rgb]{0.9,0.9,0.9}{
            \makecell[{{p{\linewidth}}}]{
                \texttt{\tiny{[P1$\langle$GM]}}
                \texttt{Es esmu jūsu spēles vadītājs, un jūs spēlējat sadarbības spēli ar šādu režģi kā spēles laukumu:} \\
\\ 
\texttt{```} \\
\texttt{ 1234567} \\
\texttt{╔═══╤═╤═╗} \\
\texttt{║◌◌I│W│◌║ 1} \\
\texttt{║◌◌◌└─┤◌║ 2} \\
\texttt{║◌◌◌TC├─╢ 3} \\
\texttt{╟─┐◌┌─┤◌║ 4} \\
\texttt{║◌│◌│◌│◌║ 5} \\
\texttt{╟─┼─┼─┼─╢ 6} \\
\texttt{║◌│◌│H│◌║ 7} \\
\texttt{╚═╧═╧═╧═╝} \\
\\ 
\texttt{```} \\
\\ 
\texttt{* Augšējā malā ir attēlotas x koordinātas, kas pieaug pa labi, un labajā malā y koordinātas, kas pieaug uz leju.} \\
\texttt{* Uz jūsu režģa ir nejauši izvietoti šādi objekti: 'W', 'I', 'T', 'C', 'H'.} \\
\\ 
\\ 
\texttt{Otrs spēlētājs redz spēles laukuma variāciju, kurā objekti ir izvietoti citās nejaušās vietās. Jūs nevarat redzēt otra spēlētāja laukumu, un viņš nevar redzēt jūsējo.} \\
\\ 
\\ 
\texttt{**Mērķis:**} \\
\\ 
\texttt{Abiem spēlētājiem ir jāpārvieto objekti uz sava fona tā, lai vienādi objekti nonāktu vienādās koordinātēs. Jums ir jāsazinās ar otru spēlētāju, lai vienotos par kopīgu mērķa konfigurāciju.} \\
\\ 
\\ 
\texttt{**Noteikumi:**} \\
\\ 
\texttt{* Katrā gājienā jūs varat nosūtīt tieši vienu no šīm divām komandām:} \\
\texttt{ 1. `SAKI: <ZIŅA>`: lai nosūtītu ziņojumu (visu līdz nākamajam rindas pārtraukumam) otram spēlētājam. Es to pārsūtīšu jūsu partnerim.} \\
\texttt{2. `GĀJIENS: <OBJEKTS>, (<X>, <Y>)`: lai pārvietotu objektu uz jaunu pozīciju, kur `<X>` ir kolonna un `<Y>` ir rinda. Es jūs informēšu, vai jūsu gājiens bija veiksmīgs vai nē.} \\
\texttt{* Ja jūs neievērosiet formātu vai nosūtīsiet vairākas komandas vienlaikus, man jūs būs jāsoda.} \\
\texttt{* Ja abi spēlētāji kopā uzkrāj vairāk nekā 12 sodus, jūs abi zaudējat spēli.} \\
\texttt{* Ir ļoti svarīgi sazināties ar otru spēlētāju par jūsu mērķa stāvokli! *Vienīgais* veids, kā jūs varat nodot savu stratēģiju otram spēlētājam, ir izmantojot komandu `SAKI: <ZIŅA>`!} \\
\\ 
\\ 
\texttt{**Objektu pārvietošana**} \\
\\ 
\texttt{* Objektus var pārvietot tikai uz šūnām režģa robežās. Mērķa šūnai jābūt tukšai, t.i., tajā jābūt tikai simbolam '◌'.} \\
\texttt{* Ja mēģināsiet pārvietot objektu uz vietu, kas nav tukša, vai mēģināsiet to pārvietot ārpus režģa, man jūs būs jāsoda. Jūs saņemat vēl vienu mēģinājumu.} \\
\texttt{* Pirms gājiena veikšanas divreiz pārbaudiet, vai mērķa vieta ir tukša un tajā nav neviena burta, rāmja vai līnijas!} \\
\\ 
\\ 
\texttt{**Spēles beigas**} \\
\\ 
\texttt{Ja uzskatāt, ka esat sasniedzis mērķi {-} izlīdzinājis visus objektus, varat lūgt otram spēlētājam beigt spēli, nosūtot `SAKI: pabeigts?`. Ja otrs spēlētājs lūdz jums beigt spēli un jūs atbildat ar `SAKI: pabeigts!`, spēle beigsies.} \\
\\ 
\texttt{Abi spēlētāji uzvar, ja spēle beidzas 20 raundos, kur viens raunds ir definēts kā divi spēlētāji, kas katrs nosūta derīgu komandu.} \\
\\ 
\texttt{**Punktu skaitīšana:**} \\
\\ 
\texttt{Jo tuvāk vienādi objekti atrodas abos spēles laukumos, jo vairāk punktu jūs saņemat. Sodi samazina jūsu punktus. Vai varat pārspēt rekordu?} \\
\\ 
\\ 
\texttt{Lūdzu, nosūtiet ziņu otram spēlētājam, lai sāktu spēli!} \\
            }
        }
    }
    & \\ \\

    \theutterance \stepcounter{utterance}  
    & \multicolumn{4}{p{0.6\linewidth}}{
        \cellcolor[rgb]{0.9,0.9,0.9}{
            \makecell[{{p{\linewidth}}}]{
                \texttt{\tiny{[P1$\rangle$GM]}}
                \texttt{Labi, sāksim spēli! Šeit ir mana pirmā ziņa, lai sāktu koordinēties ar tevi par mērķa konfigurāciju:} \\
\\ 
\texttt{```} \\
\texttt{SAKI: Sveiciens! Manā laukumā objekti izvietoti šādi:} \\
\texttt{{-} 'I' atrodas (3,1)} \\
\texttt{{-} 'W' atrodas (5,1)} \\
\texttt{{-} 'T' atrodas (4,3)} \\
\texttt{{-} 'C' atrodas (5,3)} \\
\texttt{{-} 'H' atrodas (5,7)} \\
\\ 
\texttt{Kādi ir objekti tavā laukumā? Mēs varētu sākt ar to, ka visiem objektiem mēģinām noteikt kopīgas koordinātes. Varbūt sāksim ar 'H'? Kur tas atrodas tavā laukumā?} \\
\texttt{```} \\
\\ 
\texttt{Gaidi tavu atbildi, lai turpinātu! 😊} \\
            }
        }
    }
    & & \\ \\

    \theutterance \stepcounter{utterance}  
    & & & \multicolumn{2}{p{0.3\linewidth}}{
        \cellcolor[rgb]{0.9,0.9,0.9}{
            \makecell[{{p{\linewidth}}}]{
                \texttt{\tiny{[GM$|$GM]}}
                \texttt{Labi, sāksim spēli! Šeit ir mana pirmā ziņa, lai sāktu koordinēties ar tevi par mērķa konfigurāciju:} \\
\\ 
\texttt{```} \\
\texttt{SAKI: Sveiciens! Manā laukumā objekti izvietoti šādi:} \\
\texttt{{-} 'I' atrodas (3,1)} \\
\texttt{{-} 'W' atrodas (5,1)} \\
\texttt{{-} 'T' atrodas (4,3)} \\
\texttt{{-} 'C' atrodas (5,3)} \\
\texttt{{-} 'H' atrodas (5,7)} \\
\\ 
\texttt{Kādi ir objekti tavā laukumā? Mēs varētu sākt ar to, ka visiem objektiem mēģinām noteikt kopīgas koordinātes. Varbūt sāksim ar 'H'? Kur tas atrodas tavā laukumā?} \\
\texttt{```} \\
\\ 
\texttt{Gaidi tavu atbildi, lai turpinātu! 😊} \\
            }
        }
    }
    & & \\ \\

    \theutterance \stepcounter{utterance}  
    & & & \multicolumn{2}{p{0.3\linewidth}}{
        \cellcolor[rgb]{0.9,0.9,0.9}{
            \makecell[{{p{\linewidth}}}]{
                \texttt{\tiny{[GM$|$GM]}}
                \texttt{Invalid format: head and tail are not empty} \\
\texttt{head: 'Labi, sāksim spēli! Šeit ir mana pirmā ziņa, lai sāktu koordinēties ar tevi par mērķa konfigurāciju:} \\
\\ 
\\ 
\texttt{'} \\
\texttt{tail: '} \\
\texttt{{-} 'I' atrodas (3,1)} \\
\texttt{{-} 'W' atrodas (5,1)} \\
\texttt{{-} 'T' atrodas (4,3)} \\
\texttt{{-} 'C' atrodas (5,3)} \\
\texttt{{-} 'H' atrodas (5,7)} \\
\\ 
\texttt{Kādi ir objekti tavā laukumā? Mēs varētu sākt ar to, ka visiem objektiem mēģinām noteikt kopīgas koordinātes. Varbūt sāksim ar 'H'? Kur tas atrodas tavā laukumā?} \\
\\ 
\\ 
\texttt{Gaidi tavu atbildi, lai turpinātu! 😊'} \\
            }
        }
    }
    & & \\ \\

    \theutterance \stepcounter{utterance}  
    & & \multicolumn{4}{p{0.6\linewidth}}{
        \cellcolor[rgb]{0.9,0.9,0.9}{
            \makecell[{{p{\linewidth}}}]{
                \texttt{\tiny{[P1$\langle$GM]}}
                \texttt{Sods: Jūsu ziņojumā nedrīkst būt nekas pirms vai pēc komandas!} \\
\\ 
\texttt{Pārliecinieties, ka jūsu atbildē ir tikai `SAKI: <ZIŅA>` vai `GĀJIENS: <OBJEKTS>, (<X>, <Y>)`, un nekas cits!} \\
\\ 
\texttt{Jūs kopā esat uzkrājuši 1 no 12 sodiem.} \\
\texttt{Lūdzu, mēģiniet vēlreiz!} \\
            }
        }
    }
    & \\ \\

    \theutterance \stepcounter{utterance}  
    & \multicolumn{4}{p{0.6\linewidth}}{
        \cellcolor[rgb]{0.9,0.9,0.9}{
            \makecell[{{p{\linewidth}}}]{
                \texttt{\tiny{[P1$\rangle$GM]}}
                \texttt{```} \\
\texttt{SAKI: Sveiciens! Manā laukumā 'H' atrodas (3,6). Kur atrodas tavs 'H'? Mēs varētu sākt ar to.} \\
\texttt{```} \\
            }
        }
    }
    & & \\ \\

    \theutterance \stepcounter{utterance}  
    & & & \multicolumn{2}{p{0.3\linewidth}}{
        \cellcolor[rgb]{0.9,0.9,0.9}{
            \makecell[{{p{\linewidth}}}]{
                \texttt{\tiny{[GM$|$GM]}}
                \texttt{```} \\
\texttt{SAKI: Sveiciens! Manā laukumā 'H' atrodas (3,6). Kur atrodas tavs 'H'? Mēs varētu sākt ar to.} \\
\texttt{```} \\
            }
        }
    }
    & & \\ \\

    \theutterance \stepcounter{utterance}  
    & & \multicolumn{4}{p{0.6\linewidth}}{
        \cellcolor[rgb]{0.9,0.9,0.9}{
            \makecell[{{p{\linewidth}}}]{
                \texttt{\tiny{[GM$\rangle$P2]}}
                \texttt{Es esmu jūsu spēles vadītājs, un jūs spēlējat sadarbības spēli ar šādu režģi kā spēles laukumu:} \\
\\ 
\texttt{```} \\
\texttt{ 1234567} \\
\texttt{╔═══╤═╤═╗} \\
\texttt{║◌◌◌│◌│◌║ 1} \\
\texttt{║◌W◌└─┤C║ 2} \\
\texttt{║T◌◌◌◌├─╢ 3} \\
\texttt{╟─┐◌┌─┤◌║ 4} \\
\texttt{║◌│◌│H│I║ 5} \\
\texttt{╟─┼─┼─┼─╢ 6} \\
\texttt{║◌│◌│◌│◌║ 7} \\
\texttt{╚═╧═╧═╧═╝} \\
\\ 
\texttt{```} \\
\\ 
\texttt{* Augšējā malā ir attēlotas x koordinātas, kas pieaug pa labi, un labajā malā y koordinātas, kas pieaug uz leju.} \\
\texttt{* Uz jūsu režģa ir nejauši izvietoti šādi objekti: 'W', 'I', 'T', 'C', 'H'.} \\
\\ 
\\ 
\texttt{Otrs spēlētājs redz spēles laukuma variāciju, kurā objekti ir izvietoti citās nejaušās vietās. Jūs nevarat redzēt otra spēlētāja laukumu, un viņš nevar redzēt jūsējo.} \\
\\ 
\\ 
\texttt{**Mērķis:**} \\
\\ 
\texttt{Abiem spēlētājiem ir jāpārvieto objekti uz sava fona tā, lai vienādi objekti nonāktu vienādās koordinātēs. Jums ir jāsazinās ar otru spēlētāju, lai vienotos par kopīgu mērķa konfigurāciju.} \\
\\ 
\\ 
\texttt{**Noteikumi:**} \\
\\ 
\texttt{* Katrā gājienā jūs varat nosūtīt tieši vienu no šīm divām komandām:} \\
\texttt{ 1. `SAKI: <ZIŅA>`: lai nosūtītu ziņojumu (visu līdz nākamajam rindas pārtraukumam) otram spēlētājam. Es to pārsūtīšu jūsu partnerim.} \\
\texttt{2. `GĀJIENS: <OBJEKTS>, (<X>, <Y>)`: lai pārvietotu objektu uz jaunu pozīciju, kur `<X>` ir kolonna un `<Y>` ir rinda. Es jūs informēšu, vai jūsu gājiens bija veiksmīgs vai nē.} \\
\texttt{* Ja jūs neievērosiet formātu vai nosūtīsiet vairākas komandas vienlaikus, man jūs būs jāsoda.} \\
\texttt{* Ja abi spēlētāji kopā uzkrāj vairāk nekā 12 sodus, jūs abi zaudējat spēli.} \\
\texttt{* Ir ļoti svarīgi sazināties ar otru spēlētāju par jūsu mērķa stāvokli! *Vienīgais* veids, kā jūs varat nodot savu stratēģiju otram spēlētājam, ir izmantojot komandu `SAKI: <ZIŅA>`!} \\
\\ 
\\ 
\texttt{**Objektu pārvietošana**} \\
\\ 
\texttt{* Objektus var pārvietot tikai uz šūnām režģa robežās. Mērķa šūnai jābūt tukšai, t.i., tajā jābūt tikai simbolam '◌'.} \\
\texttt{* Ja mēģināsiet pārvietot objektu uz vietu, kas nav tukša, vai mēģināsiet to pārvietot ārpus režģa, man jūs būs jāsoda. Jūs saņemat vēl vienu mēģinājumu.} \\
\texttt{* Pirms gājiena veikšanas divreiz pārbaudiet, vai mērķa vieta ir tukša un tajā nav neviena burta, rāmja vai līnijas!} \\
\\ 
\\ 
\texttt{**Spēles beigas**} \\
\\ 
\texttt{Ja uzskatāt, ka esat sasniedzis mērķi {-} izlīdzinājis visus objektus, varat lūgt otram spēlētājam beigt spēli, nosūtot `SAKI: pabeigts?`. Ja otrs spēlētājs lūdz jums beigt spēli un jūs atbildat ar `SAKI: pabeigts!`, spēle beigsies.} \\
\\ 
\texttt{Abi spēlētāji uzvar, ja spēle beidzas 20 raundos, kur viens raunds ir definēts kā divi spēlētāji, kas katrs nosūta derīgu komandu.} \\
\\ 
\texttt{**Punktu skaitīšana:**} \\
\\ 
\texttt{Jo tuvāk vienādi objekti atrodas abos spēles laukumos, jo vairāk punktu jūs saņemat. Sodi samazina jūsu punktus. Vai varat pārspēt rekordu?} \\
\\ 
\texttt{Otrs spēlētājs sāka spēli, nosūtot šo ziņojumu:} \\
\\ 
\texttt{"Sveiciens! Manā laukumā 'H' atrodas (3,6). Kur atrodas tavs 'H'? Mēs varētu sākt ar to."} \\
\\ 
\texttt{Kāda ir jūsu pirmā komanda?} \\
            }
        }
    }
    & \\ \\

    \theutterance \stepcounter{utterance}  
    & & & \multicolumn{4}{p{0.6\linewidth}}{
        \cellcolor[rgb]{0.9,0.9,0.9}{
            \makecell[{{p{\linewidth}}}]{
                \texttt{\tiny{[GM$\langle$P2]}}
                \texttt{Saskaņā ar spēles noteikumiem un otrā spēlētāja ziņojumu, manā spēles laukumā objekts 'H' atrodas pozīcijā **(6,5)** (x=6, y=5).} \\
\\ 
\texttt{Tā kā otrs spēlētājs piedāvā sākt ar 'H' objektu, es varētu atbildēt ar šādu ziņojumu, lai sāktu koordinēties par mērķa pozīcijām:} \\
\\ 
\texttt{```} \\
\texttt{SAKI: Sveiciens! Manā laukumā 'H' ir (6,5). Piedāvāju mērķa pozīciju 'H' uz (4,5), lai tas būtu tuvāk centram. Kā tev šķiet?} \\
\texttt{```} \\
\\ 
\texttt{Šī ziņa:} \\
\texttt{1. Informē par 'H' pašreizējo pozīciju manā laukumā.} \\
\texttt{2. Piedāvā kopīgu mērķa pozīciju 'H' objektam.} \\
\texttt{3. Aicina otru spēlētāju izteikt savu viedokli.} \\
\\ 
\texttt{Pēc tam, kad būsim vienojušies par mērķa pozīcijām, varēsim sākt pārvietot objektus ar komandu `GĀJIENS`.} \\
            }
        }
    }
     \\ \\

    \theutterance \stepcounter{utterance}  
    & & & \multicolumn{2}{p{0.3\linewidth}}{
        \cellcolor[rgb]{0.9,0.9,0.9}{
            \makecell[{{p{\linewidth}}}]{
                \texttt{\tiny{[GM$|$GM]}}
                \texttt{Saskaņā ar spēles noteikumiem un otrā spēlētāja ziņojumu, manā spēles laukumā objekts 'H' atrodas pozīcijā **(6,5)** (x=6, y=5).} \\
\\ 
\texttt{Tā kā otrs spēlētājs piedāvā sākt ar 'H' objektu, es varētu atbildēt ar šādu ziņojumu, lai sāktu koordinēties par mērķa pozīcijām:} \\
\\ 
\texttt{```} \\
\texttt{SAKI: Sveiciens! Manā laukumā 'H' ir (6,5). Piedāvāju mērķa pozīciju 'H' uz (4,5), lai tas būtu tuvāk centram. Kā tev šķiet?} \\
\texttt{```} \\
\\ 
\texttt{Šī ziņa:} \\
\texttt{1. Informē par 'H' pašreizējo pozīciju manā laukumā.} \\
\texttt{2. Piedāvā kopīgu mērķa pozīciju 'H' objektam.} \\
\texttt{3. Aicina otru spēlētāju izteikt savu viedokli.} \\
\\ 
\texttt{Pēc tam, kad būsim vienojušies par mērķa pozīcijām, varēsim sākt pārvietot objektus ar komandu `GĀJIENS`.} \\
            }
        }
    }
    & & \\ \\

    \theutterance \stepcounter{utterance}  
    & & & \multicolumn{2}{p{0.3\linewidth}}{
        \cellcolor[rgb]{0.9,0.9,0.9}{
            \makecell[{{p{\linewidth}}}]{
                \texttt{\tiny{[GM$|$GM]}}
                \texttt{Invalid format: head and tail are not empty} \\
\texttt{head: 'Saskaņā ar spēles noteikumiem un otrā spēlētāja ziņojumu, manā spēles laukumā objekts 'H' atrodas pozīcijā **(6,5)** (x=6, y=5).} \\
\\ 
\texttt{Tā kā otrs spēlētājs piedāvā sākt ar 'H' objektu, es varētu atbildēt ar šādu ziņojumu, lai sāktu koordinēties par mērķa pozīcijām:} \\
\\ 
\\ 
\texttt{'} \\
\texttt{tail: '} \\
\\ 
\\ 
\texttt{Šī ziņa:} \\
\texttt{1. Informē par 'H' pašreizējo pozīciju manā laukumā.} \\
\texttt{2. Piedāvā kopīgu mērķa pozīciju 'H' objektam.} \\
\texttt{3. Aicina otru spēlētāju izteikt savu viedokli.} \\
\\ 
\texttt{Pēc tam, kad būsim vienojušies par mērķa pozīcijām, varēsim sākt pārvietot objektus ar komandu GĀJIENS.'} \\
            }
        }
    }
    & & \\ \\

    \theutterance \stepcounter{utterance}  
    & & \multicolumn{4}{p{0.6\linewidth}}{
        \cellcolor[rgb]{0.9,0.9,0.9}{
            \makecell[{{p{\linewidth}}}]{
                \texttt{\tiny{[GM$\rangle$P2]}}
                \texttt{Sods: Jūsu ziņojumā nedrīkst būt nekas pirms vai pēc komandas!} \\
\\ 
\texttt{Pārliecinieties, ka jūsu atbildē ir tikai `SAKI: <ZIŅA>` vai `GĀJIENS: <OBJEKTS>, (<X>, <Y>)`, un nekas cits!} \\
\\ 
\texttt{Jūs kopā esat uzkrājuši 2 no 12 sodiem.} \\
\texttt{Lūdzu, mēģiniet vēlreiz!} \\
            }
        }
    }
    & \\ \\

    \theutterance \stepcounter{utterance}  
    & & & \multicolumn{4}{p{0.6\linewidth}}{
        \cellcolor[rgb]{0.9,0.9,0.9}{
            \makecell[{{p{\linewidth}}}]{
                \texttt{\tiny{[GM$\langle$P2]}}
                \texttt{Piedodiet par kļūdu! Šoreiz precīzi ievērošu formātu:} \\
\\ 
\texttt{```} \\
\texttt{SAKI: Man 'H' ir (6,5). Piedāvāju mērķi (4,5). Vai der?} \\
\texttt{```} \\
            }
        }
    }
     \\ \\

    \theutterance \stepcounter{utterance}  
    & & & \multicolumn{2}{p{0.3\linewidth}}{
        \cellcolor[rgb]{0.9,0.9,0.9}{
            \makecell[{{p{\linewidth}}}]{
                \texttt{\tiny{[GM$|$GM]}}
                \texttt{Piedodiet par kļūdu! Šoreiz precīzi ievērošu formātu:} \\
\\ 
\texttt{```} \\
\texttt{SAKI: Man 'H' ir (6,5). Piedāvāju mērķi (4,5). Vai der?} \\
\texttt{```} \\
            }
        }
    }
    & & \\ \\

    \theutterance \stepcounter{utterance}  
    & & & \multicolumn{2}{p{0.3\linewidth}}{
        \cellcolor[rgb]{0.9,0.9,0.9}{
            \makecell[{{p{\linewidth}}}]{
                \texttt{\tiny{[GM$|$GM]}}
                \texttt{Invalid format: head is not empty: 'Piedodiet par kļūdu! Šoreiz precīzi ievērošu formātu:} \\
\\ 
\\ 
\texttt{'} \\
            }
        }
    }
    & & \\ \\

    \theutterance \stepcounter{utterance}  
    & & \multicolumn{4}{p{0.6\linewidth}}{
        \cellcolor[rgb]{0.9,0.9,0.9}{
            \makecell[{{p{\linewidth}}}]{
                \texttt{\tiny{[GM$\rangle$P2]}}
                \texttt{Sods: Jūsu ziņojumā nedrīkst būt nekas pirms komandas!} \\
\\ 
\texttt{Pārliecinieties, ka jūsu atbildē ir tikai `SAKI: <ZIŅA>` vai `GĀJIENS: <OBJEKTS>, (<X>, <Y>)`, un nekas cits!} \\
\\ 
\texttt{Jūs kopā esat uzkrājuši 3 no 12 sodiem.} \\
\texttt{Lūdzu, mēģiniet vēlreiz!} \\
            }
        }
    }
    & \\ \\

    \theutterance \stepcounter{utterance}  
    & & & \multicolumn{4}{p{0.6\linewidth}}{
        \cellcolor[rgb]{0.9,0.9,0.9}{
            \makecell[{{p{\linewidth}}}]{
                \texttt{\tiny{[GM$\langle$P2]}}
                \texttt{```} \\
\texttt{SAKI: H (6,5). Mērķis (4,5)?} \\
\texttt{```} \\
            }
        }
    }
     \\ \\

    \theutterance \stepcounter{utterance}  
    & & & \multicolumn{2}{p{0.3\linewidth}}{
        \cellcolor[rgb]{0.9,0.9,0.9}{
            \makecell[{{p{\linewidth}}}]{
                \texttt{\tiny{[GM$|$GM]}}
                \texttt{```} \\
\texttt{SAKI: H (6,5). Mērķis (4,5)?} \\
\texttt{```} \\
            }
        }
    }
    & & \\ \\

    \theutterance \stepcounter{utterance}  
    & & \multicolumn{4}{p{0.6\linewidth}}{
        \cellcolor[rgb]{0.9,0.9,0.9}{
            \makecell[{{p{\linewidth}}}]{
                \texttt{\tiny{[P1$\langle$GM]}}
                \texttt{Jūsu ziņojums ir nodots otram spēlētājam.} \\
\\ 
\\ 
\texttt{Pašlaik jūs spēlējat kārtu 1 no maksimāli 20.} \\
\texttt{Jūs kopā esat uzkrājuši 3 no 12 sodiem.} \\
\texttt{Otrs spēlētājs nosūtīja šo ziņojumu:} \\
\\ 
\\ 
\texttt{"H (6,5). Mērķis (4,5)?"} \\
\\ 
\\ 
\texttt{Kāda ir jūsu nākamā komanda?} \\
            }
        }
    }
    & \\ \\

    \theutterance \stepcounter{utterance}  
    & \multicolumn{4}{p{0.6\linewidth}}{
        \cellcolor[rgb]{0.9,0.9,0.9}{
            \makecell[{{p{\linewidth}}}]{
                \texttt{\tiny{[P1$\rangle$GM]}}
                \texttt{```} \\
\texttt{SAKI: Labi, es pārvietoju savu 'H' uz (4,5). Tagad pārbaudīsim citu objektu. Kur atrodas tavs 'W'?} \\
\texttt{```} \\
\\ 
\texttt{*(Pēc šīs ziņas es arī veikšu gājienu, lai pārvietotu 'H' uz (4,5), ja tas ir iespējams.)*} \\
\\ 
\texttt{```} \\
\texttt{GĀJIENS: H, (4,5)} \\
\texttt{```} \\
            }
        }
    }
    & & \\ \\

    \theutterance \stepcounter{utterance}  
    & & & \multicolumn{2}{p{0.3\linewidth}}{
        \cellcolor[rgb]{0.9,0.9,0.9}{
            \makecell[{{p{\linewidth}}}]{
                \texttt{\tiny{[GM$|$GM]}}
                \texttt{```} \\
\texttt{SAKI: Labi, es pārvietoju savu 'H' uz (4,5). Tagad pārbaudīsim citu objektu. Kur atrodas tavs 'W'?} \\
\texttt{```} \\
\\ 
\texttt{*(Pēc šīs ziņas es arī veikšu gājienu, lai pārvietotu 'H' uz (4,5), ja tas ir iespējams.)*} \\
\\ 
\texttt{```} \\
\texttt{GĀJIENS: H, (4,5)} \\
\texttt{```} \\
            }
        }
    }
    & & \\ \\

    \theutterance \stepcounter{utterance}  
    & & & \multicolumn{2}{p{0.3\linewidth}}{
        \cellcolor[rgb]{0.9,0.9,0.9}{
            \makecell[{{p{\linewidth}}}]{
                \texttt{\tiny{[GM$|$GM]}}
                \texttt{Invalid response format: SAKI: Labi, es pārvietoju savu 'H' uz (4,5). Tagad pārbaudīsim citu objektu. Kur atrodas tavs 'W'?} \\
\\ 
\\ 
\texttt{*(Pēc šīs ziņas es arī veikšu gājienu, lai pārvietotu 'H' uz (4,5), ja tas ir iespējams.)*} \\
\\ 
\\ 
\texttt{GĀJIENS: H, (4,5)} \\
            }
        }
    }
    & & \\ \\

    \theutterance \stepcounter{utterance}  
    & & \multicolumn{4}{p{0.6\linewidth}}{
        \cellcolor[rgb]{0.9,0.9,0.9}{
            \makecell[{{p{\linewidth}}}]{
                \texttt{\tiny{[P1$\langle$GM]}}
                \texttt{Sods: Jūsu ziņojumā ir vairāk nekā viena komanda!} \\
\\ 
\texttt{Pārliecinieties, ka jūsu atbildē ir tikai `SAKI: <ZIŅA>` vai `GĀJIENS: <OBJEKTS>, (<X>, <Y>)`, un nekas cits!} \\
\\ 
\texttt{Jūs kopā esat uzkrājuši 4 no 12 sodiem.} \\
\texttt{Lūdzu, mēģiniet vēlreiz!} \\
            }
        }
    }
    & \\ \\

    \theutterance \stepcounter{utterance}  
    & \multicolumn{4}{p{0.6\linewidth}}{
        \cellcolor[rgb]{0.9,0.9,0.9}{
            \makecell[{{p{\linewidth}}}]{
                \texttt{\tiny{[P1$\rangle$GM]}}
                \texttt{```} \\
\texttt{GĀJIENS: H, (4,5)} \\
\texttt{```} \\
            }
        }
    }
    & & \\ \\

    \theutterance \stepcounter{utterance}  
    & & & \multicolumn{2}{p{0.3\linewidth}}{
        \cellcolor[rgb]{0.9,0.9,0.9}{
            \makecell[{{p{\linewidth}}}]{
                \texttt{\tiny{[GM$|$GM]}}
                \texttt{```} \\
\texttt{GĀJIENS: H, (4,5)} \\
\texttt{```} \\
            }
        }
    }
    & & \\ \\

    \theutterance \stepcounter{utterance}  
    & & & \multicolumn{2}{p{0.3\linewidth}}{
        \cellcolor[rgb]{0.9,0.9,0.9}{
            \makecell[{{p{\linewidth}}}]{
                \texttt{\tiny{[GM$|$GM]}}
                \texttt{Sods: (4,5) nav tukšs, bet satur '│'.} \\
\\ 
\texttt{Jūs kopā esat uzkrājuši 5 no 12 sodiem.} \\
\texttt{Lūdzu, mēģiniet vēlreiz!} \\
            }
        }
    }
    & & \\ \\

    \theutterance \stepcounter{utterance}  
    & & \multicolumn{4}{p{0.6\linewidth}}{
        \cellcolor[rgb]{0.9,0.9,0.9}{
            \makecell[{{p{\linewidth}}}]{
                \texttt{\tiny{[P1$\langle$GM]}}
                \texttt{Sods: (4,5) nav tukšs, bet satur '│'.} \\
\\ 
\texttt{Jūs kopā esat uzkrājuši 5 no 12 sodiem.} \\
\texttt{Lūdzu, mēģiniet vēlreiz!} \\
            }
        }
    }
    & \\ \\

    \theutterance \stepcounter{utterance}  
    & \multicolumn{4}{p{0.6\linewidth}}{
        \cellcolor[rgb]{0.9,0.9,0.9}{
            \makecell[{{p{\linewidth}}}]{
                \texttt{\tiny{[P1$\rangle$GM]}}
                \texttt{```} \\
\texttt{SAKI: Sapratu, (4,5) nav derīga vieta. Mēģināšu pārvietot 'H' uz (4,6). Vai tas der?} \\
\texttt{```} \\
\\ 
\texttt{*(Ja ziņa tiek pieņemta, es vēlāk nosūtīšu gājienu `GĀJIENS: H, (4,6)`)*} \\
            }
        }
    }
    & & \\ \\

    \theutterance \stepcounter{utterance}  
    & & & \multicolumn{2}{p{0.3\linewidth}}{
        \cellcolor[rgb]{0.9,0.9,0.9}{
            \makecell[{{p{\linewidth}}}]{
                \texttt{\tiny{[GM$|$GM]}}
                \texttt{```} \\
\texttt{SAKI: Sapratu, (4,5) nav derīga vieta. Mēģināšu pārvietot 'H' uz (4,6). Vai tas der?} \\
\texttt{```} \\
\\ 
\texttt{*(Ja ziņa tiek pieņemta, es vēlāk nosūtīšu gājienu `GĀJIENS: H, (4,6)`)*} \\
            }
        }
    }
    & & \\ \\

    \theutterance \stepcounter{utterance}  
    & & & \multicolumn{2}{p{0.3\linewidth}}{
        \cellcolor[rgb]{0.9,0.9,0.9}{
            \makecell[{{p{\linewidth}}}]{
                \texttt{\tiny{[GM$|$GM]}}
                \texttt{Invalid response format: SAKI: Sapratu, (4,5) nav derīga vieta. Mēģināšu pārvietot 'H' uz (4,6). Vai tas der?} \\
\\ 
\\ 
\texttt{*(Ja ziņa tiek pieņemta, es vēlāk nosūtīšu gājienu GĀJIENS: H, (4,6))*} \\
            }
        }
    }
    & & \\ \\

    \theutterance \stepcounter{utterance}  
    & & \multicolumn{4}{p{0.6\linewidth}}{
        \cellcolor[rgb]{0.9,0.9,0.9}{
            \makecell[{{p{\linewidth}}}]{
                \texttt{\tiny{[P1$\langle$GM]}}
                \texttt{Sods: Jūsu ziņojumā ir vairāk nekā viena komanda!} \\
\\ 
\texttt{Pārliecinieties, ka jūsu atbildē ir tikai `SAKI: <ZIŅA>` vai `GĀJIENS: <OBJEKTS>, (<X>, <Y>)`, un nekas cits!} \\
\\ 
\texttt{Jūs kopā esat uzkrājuši 6 no 12 sodiem.} \\
\texttt{Lūdzu, mēģiniet vēlreiz!} \\
            }
        }
    }
    & \\ \\

    \theutterance \stepcounter{utterance}  
    & \multicolumn{4}{p{0.6\linewidth}}{
        \cellcolor[rgb]{0.9,0.9,0.9}{
            \makecell[{{p{\linewidth}}}]{
                \texttt{\tiny{[P1$\rangle$GM]}}
                \texttt{```} \\
\texttt{SAKI: Mēģināšu H uz (4,6). Pārbaudīsim, vai tas der.} \\
\texttt{```} \\
            }
        }
    }
    & & \\ \\

    \theutterance \stepcounter{utterance}  
    & & & \multicolumn{2}{p{0.3\linewidth}}{
        \cellcolor[rgb]{0.9,0.9,0.9}{
            \makecell[{{p{\linewidth}}}]{
                \texttt{\tiny{[GM$|$GM]}}
                \texttt{```} \\
\texttt{SAKI: Mēģināšu H uz (4,6). Pārbaudīsim, vai tas der.} \\
\texttt{```} \\
            }
        }
    }
    & & \\ \\

    \theutterance \stepcounter{utterance}  
    & & \multicolumn{4}{p{0.6\linewidth}}{
        \cellcolor[rgb]{0.9,0.9,0.9}{
            \makecell[{{p{\linewidth}}}]{
                \texttt{\tiny{[GM$\rangle$P2]}}
                \texttt{Jūsu ziņojums ir nodots otram spēlētājam.} \\
\\ 
\\ 
\texttt{Pašlaik jūs spēlējat kārtu 2 no maksimāli 20.} \\
\texttt{Jūs kopā esat uzkrājuši 6 no 12 sodiem.} \\
\texttt{Otrs spēlētājs nosūtīja šo ziņojumu:} \\
\\ 
\\ 
\texttt{"Mēģināšu H uz (4,6). Pārbaudīsim, vai tas der."} \\
\\ 
\\ 
\texttt{Kāda ir jūsu nākamā komanda?} \\
            }
        }
    }
    & \\ \\

    \theutterance \stepcounter{utterance}  
    & & & \multicolumn{4}{p{0.6\linewidth}}{
        \cellcolor[rgb]{0.9,0.9,0.9}{
            \makecell[{{p{\linewidth}}}]{
                \texttt{\tiny{[GM$\langle$P2]}}
                \texttt{```} \\
\texttt{GĀJIENS: H, (4,5)} \\
\texttt{```} \\
\\ 
\texttt{(Es pārvietoju savu 'H' uz (4,5), lai saskaņotu ar iepriekš piedāvāto mērķi. Ja otrs spēlētājs pārvietos savu 'H' uz (4,6), mums būs jāsaskaņo vēlreiz.)} \\
            }
        }
    }
     \\ \\

    \theutterance \stepcounter{utterance}  
    & & & \multicolumn{2}{p{0.3\linewidth}}{
        \cellcolor[rgb]{0.9,0.9,0.9}{
            \makecell[{{p{\linewidth}}}]{
                \texttt{\tiny{[GM$|$GM]}}
                \texttt{```} \\
\texttt{GĀJIENS: H, (4,5)} \\
\texttt{```} \\
\\ 
\texttt{(Es pārvietoju savu 'H' uz (4,5), lai saskaņotu ar iepriekš piedāvāto mērķi. Ja otrs spēlētājs pārvietos savu 'H' uz (4,6), mums būs jāsaskaņo vēlreiz.)} \\
            }
        }
    }
    & & \\ \\

    \theutterance \stepcounter{utterance}  
    & & & \multicolumn{2}{p{0.3\linewidth}}{
        \cellcolor[rgb]{0.9,0.9,0.9}{
            \makecell[{{p{\linewidth}}}]{
                \texttt{\tiny{[GM$|$GM]}}
                \texttt{Invalid format: tail is not empty: '} \\
\\ 
\\ 
\texttt{(Es pārvietoju savu 'H' uz (4,5), lai saskaņotu ar iepriekš piedāvāto mērķi. Ja otrs spēlētājs pārvietos savu 'H' uz (4,6), mums būs jāsaskaņo vēlreiz.)} \\
            }
        }
    }
    & & \\ \\

    \theutterance \stepcounter{utterance}  
    & & \multicolumn{4}{p{0.6\linewidth}}{
        \cellcolor[rgb]{0.9,0.9,0.9}{
            \makecell[{{p{\linewidth}}}]{
                \texttt{\tiny{[GM$\rangle$P2]}}
                \texttt{Sods: Jūsu ziņojumā nedrīkst būt nekas pēc komandas!} \\
\\ 
\texttt{Pārliecinieties, ka jūsu atbildē ir tikai `SAKI: <ZIŅA>` vai `GĀJIENS: <OBJEKTS>, (<X>, <Y>)`, un nekas cits!} \\
\\ 
\texttt{Jūs kopā esat uzkrājuši 7 no 12 sodiem.} \\
\texttt{Lūdzu, mēģiniet vēlreiz!} \\
            }
        }
    }
    & \\ \\

    \theutterance \stepcounter{utterance}  
    & & & \multicolumn{4}{p{0.6\linewidth}}{
        \cellcolor[rgb]{0.9,0.9,0.9}{
            \makecell[{{p{\linewidth}}}]{
                \texttt{\tiny{[GM$\langle$P2]}}
                \texttt{```} \\
\texttt{GĀJIENS: H, (4,5)} \\
\texttt{```} \\
            }
        }
    }
     \\ \\

    \theutterance \stepcounter{utterance}  
    & & & \multicolumn{2}{p{0.3\linewidth}}{
        \cellcolor[rgb]{0.9,0.9,0.9}{
            \makecell[{{p{\linewidth}}}]{
                \texttt{\tiny{[GM$|$GM]}}
                \texttt{```} \\
\texttt{GĀJIENS: H, (4,5)} \\
\texttt{```} \\
            }
        }
    }
    & & \\ \\

    \theutterance \stepcounter{utterance}  
    & & & \multicolumn{2}{p{0.3\linewidth}}{
        \cellcolor[rgb]{0.9,0.9,0.9}{
            \makecell[{{p{\linewidth}}}]{
                \texttt{\tiny{[GM$|$GM]}}
                \texttt{Sods: (4,5) nav tukšs, bet satur '│'.} \\
\\ 
\texttt{Jūs kopā esat uzkrājuši 8 no 12 sodiem.} \\
\texttt{Lūdzu, mēģiniet vēlreiz!} \\
            }
        }
    }
    & & \\ \\

    \theutterance \stepcounter{utterance}  
    & & \multicolumn{4}{p{0.6\linewidth}}{
        \cellcolor[rgb]{0.9,0.9,0.9}{
            \makecell[{{p{\linewidth}}}]{
                \texttt{\tiny{[GM$\rangle$P2]}}
                \texttt{Sods: (4,5) nav tukšs, bet satur '│'.} \\
\\ 
\texttt{Jūs kopā esat uzkrājuši 8 no 12 sodiem.} \\
\texttt{Lūdzu, mēģiniet vēlreiz!} \\
            }
        }
    }
    & \\ \\

    \theutterance \stepcounter{utterance}  
    & & & \multicolumn{4}{p{0.6\linewidth}}{
        \cellcolor[rgb]{0.9,0.9,0.9}{
            \makecell[{{p{\linewidth}}}]{
                \texttt{\tiny{[GM$\langle$P2]}}
                \texttt{```} \\
\texttt{SAKI: (4,5) nav tukšs. Piedāvāju H mērķi (5,5). Vai der?} \\
\texttt{```} \\
\\ 
\texttt{(Pārbaudīju režģi: (4,5) satur 'I', bet (5,5) ir tukša '◌'. Piedāvāju jaunu mērķi.)} \\
            }
        }
    }
     \\ \\

    \theutterance \stepcounter{utterance}  
    & & & \multicolumn{2}{p{0.3\linewidth}}{
        \cellcolor[rgb]{0.9,0.9,0.9}{
            \makecell[{{p{\linewidth}}}]{
                \texttt{\tiny{[GM$|$GM]}}
                \texttt{```} \\
\texttt{SAKI: (4,5) nav tukšs. Piedāvāju H mērķi (5,5). Vai der?} \\
\texttt{```} \\
\\ 
\texttt{(Pārbaudīju režģi: (4,5) satur 'I', bet (5,5) ir tukša '◌'. Piedāvāju jaunu mērķi.)} \\
            }
        }
    }
    & & \\ \\

    \theutterance \stepcounter{utterance}  
    & & & \multicolumn{2}{p{0.3\linewidth}}{
        \cellcolor[rgb]{0.9,0.9,0.9}{
            \makecell[{{p{\linewidth}}}]{
                \texttt{\tiny{[GM$|$GM]}}
                \texttt{Invalid format: tail is not empty: '} \\
\\ 
\\ 
\texttt{(Pārbaudīju režģi: (4,5) satur 'I', bet (5,5) ir tukša '◌'. Piedāvāju jaunu mērķi.)} \\
            }
        }
    }
    & & \\ \\

    \theutterance \stepcounter{utterance}  
    & & \multicolumn{4}{p{0.6\linewidth}}{
        \cellcolor[rgb]{0.9,0.9,0.9}{
            \makecell[{{p{\linewidth}}}]{
                \texttt{\tiny{[GM$\rangle$P2]}}
                \texttt{Sods: Jūsu ziņojumā nedrīkst būt nekas pēc komandas!} \\
\\ 
\texttt{Pārliecinieties, ka jūsu atbildē ir tikai `SAKI: <ZIŅA>` vai `GĀJIENS: <OBJEKTS>, (<X>, <Y>)`, un nekas cits!} \\
\\ 
\texttt{Jūs kopā esat uzkrājuši 9 no 12 sodiem.} \\
\texttt{Lūdzu, mēģiniet vēlreiz!} \\
            }
        }
    }
    & \\ \\

    \theutterance \stepcounter{utterance}  
    & & & \multicolumn{4}{p{0.6\linewidth}}{
        \cellcolor[rgb]{0.9,0.9,0.9}{
            \makecell[{{p{\linewidth}}}]{
                \texttt{\tiny{[GM$\langle$P2]}}
                \texttt{```} \\
\texttt{SAKI: H mērķis (5,5)?} \\
\texttt{```} \\
            }
        }
    }
     \\ \\

    \theutterance \stepcounter{utterance}  
    & & & \multicolumn{2}{p{0.3\linewidth}}{
        \cellcolor[rgb]{0.9,0.9,0.9}{
            \makecell[{{p{\linewidth}}}]{
                \texttt{\tiny{[GM$|$GM]}}
                \texttt{```} \\
\texttt{SAKI: H mērķis (5,5)?} \\
\texttt{```} \\
            }
        }
    }
    & & \\ \\

    \theutterance \stepcounter{utterance}  
    & & \multicolumn{4}{p{0.6\linewidth}}{
        \cellcolor[rgb]{0.9,0.9,0.9}{
            \makecell[{{p{\linewidth}}}]{
                \texttt{\tiny{[P1$\langle$GM]}}
                \texttt{Jūsu ziņojums ir nodots otram spēlētājam.} \\
\\ 
\\ 
\texttt{Pašlaik jūs spēlējat kārtu 2 no maksimāli 20.} \\
\texttt{Jūs kopā esat uzkrājuši 9 no 12 sodiem.} \\
\texttt{Otrs spēlētājs nosūtīja šo ziņojumu:} \\
\\ 
\\ 
\texttt{"H mērķis (5,5)?"} \\
\\ 
\\ 
\texttt{Kāda ir jūsu nākamā komanda?} \\
            }
        }
    }
    & \\ \\

    \theutterance \stepcounter{utterance}  
    & \multicolumn{4}{p{0.6\linewidth}}{
        \cellcolor[rgb]{0.9,0.9,0.9}{
            \makecell[{{p{\linewidth}}}]{
                \texttt{\tiny{[P1$\rangle$GM]}}
                \texttt{```} \\
\texttt{GĀJIENS: H, (5,5)} \\
\texttt{```} \\
            }
        }
    }
    & & \\ \\

    \theutterance \stepcounter{utterance}  
    & & & \multicolumn{2}{p{0.3\linewidth}}{
        \cellcolor[rgb]{0.9,0.9,0.9}{
            \makecell[{{p{\linewidth}}}]{
                \texttt{\tiny{[GM$|$GM]}}
                \texttt{```} \\
\texttt{GĀJIENS: H, (5,5)} \\
\texttt{```} \\
            }
        }
    }
    & & \\ \\

    \theutterance \stepcounter{utterance}  
    & & & \multicolumn{2}{p{0.3\linewidth}}{
        \cellcolor[rgb]{0.9,0.9,0.9}{
            \makecell[{{p{\linewidth}}}]{
                \texttt{\tiny{[GM$|$GM]}}
                \texttt{H moved to (5, 5)} \\
            }
        }
    }
    & & \\ \\

    \theutterance \stepcounter{utterance}  
    & & \multicolumn{4}{p{0.6\linewidth}}{
        \cellcolor[rgb]{0.9,0.9,0.9}{
            \makecell[{{p{\linewidth}}}]{
                \texttt{\tiny{[GM$\rangle$P2]}}
                \texttt{Jūsu ziņojums ir nodots otram spēlētājam.} \\
\\ 
\\ 
\texttt{Pašlaik jūs spēlējat kārtu 3 no maksimāli 20.} \\
\texttt{Jūs kopā esat uzkrājuši 9 no 12 sodiem.} \\
\texttt{Otrs spēlētājs pārvietoja objektu savā režģī.} \\
\\ 
\\ 
\texttt{Kāda ir jūsu nākamā komanda?} \\
            }
        }
    }
    & \\ \\

    \theutterance \stepcounter{utterance}  
    & & & \multicolumn{4}{p{0.6\linewidth}}{
        \cellcolor[rgb]{0.9,0.9,0.9}{
            \makecell[{{p{\linewidth}}}]{
                \texttt{\tiny{[GM$\langle$P2]}}
                \texttt{```} \\
\texttt{GĀJIENS: H, (5,5)} \\
\texttt{```} \\
            }
        }
    }
     \\ \\

    \theutterance \stepcounter{utterance}  
    & & & \multicolumn{2}{p{0.3\linewidth}}{
        \cellcolor[rgb]{0.9,0.9,0.9}{
            \makecell[{{p{\linewidth}}}]{
                \texttt{\tiny{[GM$|$GM]}}
                \texttt{```} \\
\texttt{GĀJIENS: H, (5,5)} \\
\texttt{```} \\
            }
        }
    }
    & & \\ \\

    \theutterance \stepcounter{utterance}  
    & & & \multicolumn{2}{p{0.3\linewidth}}{
        \cellcolor[rgb]{0.9,0.9,0.9}{
            \makecell[{{p{\linewidth}}}]{
                \texttt{\tiny{[GM$|$GM]}}
                \texttt{Sods: (5,5) nav tukšs, bet satur 'H'.} \\
\\ 
\texttt{Jūs kopā esat uzkrājuši 10 no 12 sodiem.} \\
\texttt{Lūdzu, mēģiniet vēlreiz!} \\
            }
        }
    }
    & & \\ \\

    \theutterance \stepcounter{utterance}  
    & & \multicolumn{4}{p{0.6\linewidth}}{
        \cellcolor[rgb]{0.9,0.9,0.9}{
            \makecell[{{p{\linewidth}}}]{
                \texttt{\tiny{[GM$\rangle$P2]}}
                \texttt{Sods: (5,5) nav tukšs, bet satur 'H'.} \\
\\ 
\texttt{Jūs kopā esat uzkrājuši 10 no 12 sodiem.} \\
\texttt{Lūdzu, mēģiniet vēlreiz!} \\
            }
        }
    }
    & \\ \\

    \theutterance \stepcounter{utterance}  
    & & & \multicolumn{4}{p{0.6\linewidth}}{
        \cellcolor[rgb]{0.9,0.9,0.9}{
            \makecell[{{p{\linewidth}}}]{
                \texttt{\tiny{[GM$\langle$P2]}}
                \texttt{```} \\
\texttt{SAKI: H jau (5,5). Kur tavs H?} \\
\texttt{```} \\
\\ 
\texttt{(Tagad manas 'H' jau ir (5,5), bet es jautāju, kur atrodas otra spēlētāja 'H', lai saskaņotu pozīcijas. Sods par iepriekšējo kļūdu bija par to, ka (5,5) jau bija aizņemts ar 'H' {-} tāpēc es precizēju situāciju.)} \\
            }
        }
    }
     \\ \\

    \theutterance \stepcounter{utterance}  
    & & & \multicolumn{2}{p{0.3\linewidth}}{
        \cellcolor[rgb]{0.9,0.9,0.9}{
            \makecell[{{p{\linewidth}}}]{
                \texttt{\tiny{[GM$|$GM]}}
                \texttt{```} \\
\texttt{SAKI: H jau (5,5). Kur tavs H?} \\
\texttt{```} \\
\\ 
\texttt{(Tagad manas 'H' jau ir (5,5), bet es jautāju, kur atrodas otra spēlētāja 'H', lai saskaņotu pozīcijas. Sods par iepriekšējo kļūdu bija par to, ka (5,5) jau bija aizņemts ar 'H' {-} tāpēc es precizēju situāciju.)} \\
            }
        }
    }
    & & \\ \\

    \theutterance \stepcounter{utterance}  
    & & & \multicolumn{2}{p{0.3\linewidth}}{
        \cellcolor[rgb]{0.9,0.9,0.9}{
            \makecell[{{p{\linewidth}}}]{
                \texttt{\tiny{[GM$|$GM]}}
                \texttt{Invalid format: tail is not empty: '} \\
\\ 
\\ 
\texttt{(Tagad manas 'H' jau ir (5,5), bet es jautāju, kur atrodas otra spēlētāja 'H', lai saskaņotu pozīcijas. Sods par iepriekšējo kļūdu bija par to, ka (5,5) jau bija aizņemts ar 'H' {-} tāpēc es precizēju situāciju.)} \\
            }
        }
    }
    & & \\ \\

    \theutterance \stepcounter{utterance}  
    & & \multicolumn{4}{p{0.6\linewidth}}{
        \cellcolor[rgb]{0.9,0.9,0.9}{
            \makecell[{{p{\linewidth}}}]{
                \texttt{\tiny{[GM$\rangle$P2]}}
                \texttt{Sods: Jūsu ziņojumā nedrīkst būt nekas pēc komandas!} \\
\\ 
\texttt{Pārliecinieties, ka jūsu atbildē ir tikai `SAKI: <ZIŅA>` vai `GĀJIENS: <OBJEKTS>, (<X>, <Y>)`, un nekas cits!} \\
\\ 
\texttt{Jūs kopā esat uzkrājuši 11 no 12 sodiem.} \\
\texttt{Lūdzu, mēģiniet vēlreiz!} \\
            }
        }
    }
    & \\ \\

    \theutterance \stepcounter{utterance}  
    & & & \multicolumn{4}{p{0.6\linewidth}}{
        \cellcolor[rgb]{0.9,0.9,0.9}{
            \makecell[{{p{\linewidth}}}]{
                \texttt{\tiny{[GM$\langle$P2]}}
                \texttt{```} \\
\texttt{SAKI: Kur tavs H?} \\
\texttt{```} \\
            }
        }
    }
     \\ \\

    \theutterance \stepcounter{utterance}  
    & & & \multicolumn{2}{p{0.3\linewidth}}{
        \cellcolor[rgb]{0.9,0.9,0.9}{
            \makecell[{{p{\linewidth}}}]{
                \texttt{\tiny{[GM$|$GM]}}
                \texttt{```} \\
\texttt{SAKI: Kur tavs H?} \\
\texttt{```} \\
            }
        }
    }
    & & \\ \\

    \theutterance \stepcounter{utterance}  
    & & \multicolumn{4}{p{0.6\linewidth}}{
        \cellcolor[rgb]{0.9,0.9,0.9}{
            \makecell[{{p{\linewidth}}}]{
                \texttt{\tiny{[P1$\langle$GM]}}
                \texttt{Veiksmīgi pārvietots 'H' uz (5,5). Jūsu atjauninātais režģis izskatās šādi:} \\
\texttt{```} \\
\texttt{ 1234567} \\
\texttt{╔═══╤═╤═╗} \\
\texttt{║◌◌I│W│◌║ 1} \\
\texttt{║◌◌◌└─┤◌║ 2} \\
\texttt{║◌◌◌TC├─╢ 3} \\
\texttt{╟─┐◌┌─┤◌║ 4} \\
\texttt{║◌│◌│H│◌║ 5} \\
\texttt{╟─┼─┼─┼─╢ 6} \\
\texttt{║◌│◌│◌│◌║ 7} \\
\texttt{╚═╧═╧═╧═╝} \\
\texttt{```} \\
\texttt{Pašlaik jūs spēlējat kārtu 3 no maksimāli 20.} \\
\texttt{Jūs kopā esat uzkrājuši 11 no 12 sodiem.} \\
\texttt{Otrs spēlētājs nosūtīja šo ziņojumu:} \\
\\ 
\\ 
\texttt{"Kur tavs H?"} \\
\\ 
\\ 
\texttt{Kāda ir jūsu nākamā komanda?} \\
            }
        }
    }
    & \\ \\

    \theutterance \stepcounter{utterance}  
    & \multicolumn{4}{p{0.6\linewidth}}{
        \cellcolor[rgb]{0.9,0.9,0.9}{
            \makecell[{{p{\linewidth}}}]{
                \texttt{\tiny{[P1$\rangle$GM]}}
                \texttt{```} \\
\texttt{SAKI: Mans 'H' tagad ir (5,5). Kur atrodas tavs 'W'? Es redzu savu 'W' (5,1). Mēģināsim to?} \\
\texttt{```} \\
\\ 
\texttt{*(Piezīme: Mums ir tikai 1 soda punkts atlicis, tāpēc būsim uzmanīgi!)*} \\
            }
        }
    }
    & & \\ \\

    \theutterance \stepcounter{utterance}  
    & & & \multicolumn{2}{p{0.3\linewidth}}{
        \cellcolor[rgb]{0.9,0.9,0.9}{
            \makecell[{{p{\linewidth}}}]{
                \texttt{\tiny{[GM$|$GM]}}
                \texttt{```} \\
\texttt{SAKI: Mans 'H' tagad ir (5,5). Kur atrodas tavs 'W'? Es redzu savu 'W' (5,1). Mēģināsim to?} \\
\texttt{```} \\
\\ 
\texttt{*(Piezīme: Mums ir tikai 1 soda punkts atlicis, tāpēc būsim uzmanīgi!)*} \\
            }
        }
    }
    & & \\ \\

    \theutterance \stepcounter{utterance}  
    & & & \multicolumn{2}{p{0.3\linewidth}}{
        \cellcolor[rgb]{0.9,0.9,0.9}{
            \makecell[{{p{\linewidth}}}]{
                \texttt{\tiny{[GM$|$GM]}}
                \texttt{Invalid format: tail is not empty: '} \\
\\ 
\\ 
\texttt{*(Piezīme: Mums ir tikai 1 soda punkts atlicis, tāpēc būsim uzmanīgi!)*} \\
            }
        }
    }
    & & \\ \\

    \theutterance \stepcounter{utterance}  
    & & \multicolumn{4}{p{0.6\linewidth}}{
        \cellcolor[rgb]{0.9,0.9,0.9}{
            \makecell[{{p{\linewidth}}}]{
                \texttt{\tiny{[P1$\langle$GM]}}
                \texttt{Sods: Jūsu ziņojumā nedrīkst būt nekas pēc komandas!} \\
\\ 
\texttt{Pārliecinieties, ka jūsu atbildē ir tikai `SAKI: <ZIŅA>` vai `GĀJIENS: <OBJEKTS>, (<X>, <Y>)`, un nekas cits!} \\
\\ 
\texttt{Jūs kopā esat uzkrājuši 12 no 12 sodiem.} \\
\texttt{Lūdzu, mēģiniet vēlreiz!} \\
            }
        }
    }
    & \\ \\

    \theutterance \stepcounter{utterance}  
    & \multicolumn{4}{p{0.6\linewidth}}{
        \cellcolor[rgb]{0.9,0.9,0.9}{
            \makecell[{{p{\linewidth}}}]{
                \texttt{\tiny{[P1$\rangle$GM]}}
                \texttt{```} \\
\texttt{SAKI: Mans H (5,5). Tavs W?} \\
\texttt{```} \\
            }
        }
    }
    & & \\ \\

    \theutterance \stepcounter{utterance}  
    & & & \multicolumn{2}{p{0.3\linewidth}}{
        \cellcolor[rgb]{0.9,0.9,0.9}{
            \makecell[{{p{\linewidth}}}]{
                \texttt{\tiny{[GM$|$GM]}}
                \texttt{```} \\
\texttt{SAKI: Mans H (5,5). Tavs W?} \\
\texttt{```} \\
            }
        }
    }
    & & \\ \\

    \theutterance \stepcounter{utterance}  
    & & \multicolumn{4}{p{0.6\linewidth}}{
        \cellcolor[rgb]{0.9,0.9,0.9}{
            \makecell[{{p{\linewidth}}}]{
                \texttt{\tiny{[GM$\rangle$P2]}}
                \texttt{Jūsu ziņojums ir nodots otram spēlētājam.} \\
\\ 
\\ 
\texttt{Pašlaik jūs spēlējat kārtu 4 no maksimāli 20.} \\
\texttt{Jūs kopā esat uzkrājuši 12 no 12 sodiem.} \\
\texttt{Otrs spēlētājs nosūtīja šo ziņojumu:} \\
\\ 
\\ 
\texttt{"Mans H (5,5). Tavs W?"} \\
\\ 
\\ 
\texttt{Kāda ir jūsu nākamā komanda?} \\
            }
        }
    }
    & \\ \\

    \theutterance \stepcounter{utterance}  
    & & & \multicolumn{4}{p{0.6\linewidth}}{
        \cellcolor[rgb]{0.9,0.9,0.9}{
            \makecell[{{p{\linewidth}}}]{
                \texttt{\tiny{[GM$\langle$P2]}}
                \texttt{```} \\
\texttt{SAKI: Mans W (2,2). Mērķis (3,3)?} \\
\texttt{```} \\
\\ 
\texttt{(Tagad, kad abu 'H' ir (5,5), es piedāvāju sākt koordinēt 'W' objektu. Manā laukumā 'W' atrodas (2,2), un (3,3) ir tukša šūna, ko varētu izmantot kā mērķi.)} \\
            }
        }
    }
     \\ \\

    \theutterance \stepcounter{utterance}  
    & & & \multicolumn{2}{p{0.3\linewidth}}{
        \cellcolor[rgb]{0.9,0.9,0.9}{
            \makecell[{{p{\linewidth}}}]{
                \texttt{\tiny{[GM$|$GM]}}
                \texttt{```} \\
\texttt{SAKI: Mans W (2,2). Mērķis (3,3)?} \\
\texttt{```} \\
\\ 
\texttt{(Tagad, kad abu 'H' ir (5,5), es piedāvāju sākt koordinēt 'W' objektu. Manā laukumā 'W' atrodas (2,2), un (3,3) ir tukša šūna, ko varētu izmantot kā mērķi.)} \\
            }
        }
    }
    & & \\ \\

    \theutterance \stepcounter{utterance}  
    & & & \multicolumn{2}{p{0.3\linewidth}}{
        \cellcolor[rgb]{0.9,0.9,0.9}{
            \makecell[{{p{\linewidth}}}]{
                \texttt{\tiny{[GM$|$GM]}}
                \texttt{Invalid format: tail is not empty: '} \\
\\ 
\\ 
\texttt{(Tagad, kad abu 'H' ir (5,5), es piedāvāju sākt koordinēt 'W' objektu. Manā laukumā 'W' atrodas (2,2), un (3,3) ir tukša šūna, ko varētu izmantot kā mērķi.)} \\
            }
        }
    }
    & & \\ \\

    \theutterance \stepcounter{utterance}  
    & & & \multicolumn{2}{p{0.3\linewidth}}{
        \cellcolor[rgb]{0.9,0.9,0.9}{
            \makecell[{{p{\linewidth}}}]{
                \texttt{\tiny{[GM$|$GM]}}
                \texttt{Maximum number of penalties exceeded} \\
            }
        }
    }
    & & \\ \\

    \theutterance \stepcounter{utterance}  
    & & & \multicolumn{2}{p{0.3\linewidth}}{
        \cellcolor[rgb]{0.9,0.9,0.9}{
            \makecell[{{p{\linewidth}}}]{
                \texttt{\tiny{[GM$|$GM]}}
                \texttt{Initial states:} \\
\texttt{```} \\
\texttt{Player 1:} \\
\texttt{ 1234567} \\
\texttt{╔═══╤═╤═╗} \\
\texttt{║◌◌I│W│◌║ 1} \\
\texttt{║◌◌◌└─┤◌║ 2} \\
\texttt{║◌◌◌TC├─╢ 3} \\
\texttt{╟─┐◌┌─┤◌║ 4} \\
\texttt{║◌│◌│◌│◌║ 5} \\
\texttt{╟─┼─┼─┼─╢ 6} \\
\texttt{║◌│◌│H│◌║ 7} \\
\texttt{╚═╧═╧═╧═╝} \\
\\ 
\\ 
\texttt{Player 2:} \\
\texttt{ 1234567} \\
\texttt{╔═══╤═╤═╗} \\
\texttt{║◌◌◌│◌│◌║ 1} \\
\texttt{║◌W◌└─┤C║ 2} \\
\texttt{║T◌◌◌◌├─╢ 3} \\
\texttt{╟─┐◌┌─┤◌║ 4} \\
\texttt{║◌│◌│H│I║ 5} \\
\texttt{╟─┼─┼─┼─╢ 6} \\
\texttt{║◌│◌│◌│◌║ 7} \\
\texttt{╚═╧═╧═╧═╝} \\
\\ 
\texttt{```} \\
            }
        }
    }
    & & \\ \\

    \theutterance \stepcounter{utterance}  
    & & & \multicolumn{2}{p{0.3\linewidth}}{
        \cellcolor[rgb]{0.9,0.9,0.9}{
            \makecell[{{p{\linewidth}}}]{
                \texttt{\tiny{[GM$|$GM]}}
                \texttt{End states:} \\
\texttt{```} \\
\texttt{Player 1:} \\
\texttt{ 1234567} \\
\texttt{╔═══╤═╤═╗} \\
\texttt{║◌◌I│W│◌║ 1} \\
\texttt{║◌◌◌└─┤◌║ 2} \\
\texttt{║◌◌◌TC├─╢ 3} \\
\texttt{╟─┐◌┌─┤◌║ 4} \\
\texttt{║◌│◌│H│◌║ 5} \\
\texttt{╟─┼─┼─┼─╢ 6} \\
\texttt{║◌│◌│◌│◌║ 7} \\
\texttt{╚═╧═╧═╧═╝} \\
\\ 
\\ 
\texttt{Player 2:} \\
\texttt{ 1234567} \\
\texttt{╔═══╤═╤═╗} \\
\texttt{║◌◌◌│◌│◌║ 1} \\
\texttt{║◌W◌└─┤C║ 2} \\
\texttt{║T◌◌◌◌├─╢ 3} \\
\texttt{╟─┐◌┌─┤◌║ 4} \\
\texttt{║◌│◌│H│I║ 5} \\
\texttt{╟─┼─┼─┼─╢ 6} \\
\texttt{║◌│◌│◌│◌║ 7} \\
\texttt{╚═╧═╧═╧═╝} \\
\\ 
\texttt{```} \\
            }
        }
    }
    & & \\ \\

    \theutterance \stepcounter{utterance}  
    & & & \multicolumn{2}{p{0.3\linewidth}}{
        \cellcolor[rgb]{0.9,0.9,0.9}{
            \makecell[{{p{\linewidth}}}]{
                \texttt{\tiny{[GM$|$GM]}}
                \texttt{* success: False} \\
\texttt{* lose: False} \\
\texttt{* aborted: True} \\
\texttt{{-}{-}{-}{-}{-}{-}{-}} \\
\texttt{* Shifts: 0.00} \\
\texttt{* Max Shifts: 8.00} \\
\texttt{* Min Shifts: 4.00} \\
\texttt{* End Distance Sum: 14.06} \\
\texttt{* Init Distance Sum: 16.06} \\
\texttt{* Expected Distance Sum: 20.95} \\
\texttt{* Penalties: 13.00} \\
\texttt{* Max Penalties: 12.00} \\
\texttt{* Rounds: 3.00} \\
\texttt{* Max Rounds: 20.00} \\
\texttt{* Object Count: 5.00} \\
            }
        }
    }
    & & \\ \\

\end{supertabular}
}

\end{document}
